\documentclass[final]{beamer}
\usepackage[size=a0,orientation=landscape,scale=1.15]{beamerposter}
\usetheme{Madrid}
\usecolortheme{whale}
\usepackage[T1]{fontenc}
\usepackage[utf8]{inputenc}
\usepackage{lmodern,amsmath,amssymb,booktabs,array,graphicx,pgfplots}
\pgfplotsset{compat=1.18}
\definecolor{navy}{HTML}{204C70}
\definecolor{burnt}{HTML}{B85932}
\definecolor{teal}{HTML}{28786B}
\definecolor{lightteal}{HTML}{A3C6BC}
\setbeamertemplate{navigation symbols}{}
\setbeamertemplate{footline}{}
\setbeamertemplate{headline}{}
\setbeamertemplate{itemize item}{\color{navy}\textbullet}
\setbeamercolor{block title}{fg=white,bg=navy}
\setbeamercolor{block body}{fg=black,bg=white}
\setbeamerfont{block title}{size=\large,series=\bfseries}
\setbeamerfont{block body}{size=\normalsize}
\newcommand{\plotheight}{18cm}
\newcommand{\barwidth}{65pt}
\newcommand{\pointsize}{4pt}
\setbeamersize{text margin left=1.5cm,text margin right=1.5cm}
\newcommand{\smallsource}[1]{\par\medskip{\footnotesize\color{gray}#1\par}}
\setlength{\tabcolsep}{10pt}
\renewcommand{\arraystretch}{1.2}
\begin{document}
\begin{frame}[t]
\begin{center}
{\fontsize{64}{72}\selectfont\color{navy}\bfseries Custom Memory Allocator Profiling\par}
\vspace{.4cm}
{\Large Workload-specific improvements across code paths and environments\par}
\vspace{.6cm}
{\large Jacob Nance \quad Fernanda Gonzalez \quad Francisco Chagua \quad Nicholas Hooper \quad Nicholas Johnson \quad Mohamed Konsowa\par}
\vspace{.4cm}
{\normalsize CSE-625, Group 2 \qquad Fall 2026 \qquad Updated September 22, 2026\par}
\end{center}
\vspace{.6cm}
\begin{columns}[t,totalwidth=.96\paperwidth]
\begin{column}{.30\paperwidth}
\begin{block}{A representative path through a large design space}
Exhaustive profiling would cross versions, request sizes, lifetimes, concurrency, hardware, operating systems, and machine state. Even a simplified $8\times4\times3\times2\times5$ design requires \textbf{960 runs per environment}.
\medskip
We chose workloads that expose different allocator costs:
\medskip
\small
\begin{tabular}{p{.31\linewidth}p{.61\linewidth}}
\toprule Workload & Mechanism tested \\\midrule
Matrix array & Repeated row requests and placement \\
Uniform nodes & Same-size reuse with object churn \\
Nested stress & Mixed sizes, fragmentation, coalescing \\
Linked list & Tiny objects, metadata, frequent allocation \\\bottomrule
\end{tabular}\par
\normalsize
\medskip
Matched Regular/Overloaded programs use the same workload, arguments, and deterministic input within each comparison. GNU time records elapsed time and resource metrics. Checksums check output consistency.
\smallsource{Report IV and VI. The 960-run count illustrates a complete design, not completed coverage.}
\end{block}
\vspace{.6cm}
\begin{block}{What v8 changes in the code}
The allocator redirects global \texttt{new/delete} into a fixed pool. Search, splitting, alignment, metadata, and merge policy affect the cost of each request.
\medskip
V8 specifically targets separately allocated matrix rows:
\begin{itemize}
\item Split small/large regions with a reserved guard.
\item Preserve exact split sizes for packed rows.
\item Align requests of at least 2048 bytes to 64 bytes.
\item Bypass the small-object thread cache and retain deferred coalescing.
\end{itemize}
The report’s v8 linked-list runs did not finish. This is a specialization tradeoff, not a general allocator improvement.
\smallsource{Report III-J/K. Implementation descriptions support hypotheses; the study does not measure a per-function runtime breakdown.}
\end{block}
\vspace{.6cm}
\begin{block}{Statistical evidence beyond averages}
\textbf{Paired ANOVA} tests mean runtime differences, with $d_i=O_i-R_i$:
\[
F=\frac{n\bar d^2}{s_d^2},\qquad df=(1,n-1).
\]
\textbf{Directional chi-square} counts increases and decreases:
\[
\chi^2=\frac{(I-D)^2}{I+D}.
\]
For v8, expected counts are below five, so significance uses an \textbf{exact two-sided binomial sign test}. ANOVA and chi-square are separate analyses.
\medskip
The pthread report also uses Shapiro--Wilk, Levene, one-way ANOVA, and Tukey HSD. V8 environment comparisons use Welch ANOVA on paired percentage changes.
\par\medskip
V8 pairing follows run order. Paired ANOVA assumes valid pairs, independent observations across pairs, and approximately normal differences.
\smallsource{Report VII/XII and updated analysis. Threshold $\alpha=0.05$; v8 p values are unadjusted and exploratory.}
\end{block}
\end{column}
\begin{column}{.30\paperwidth}
\begin{block}{V8 runtime gains vary across groups}
\centering\renewcommand{\plotheight}{18cm}\begin{tikzpicture}\begin{axis}[width=.94\linewidth,height=\plotheight,axis lines=left,tick label style={font=\footnotesize},label style={font=\small},legend style={font=\footnotesize,draw=none},xmin=-.6,xmax=2.6,ymin=-20,ymax=60,xtick={0,1,2},xticklabels={Short,Medium,Long},ylabel={Paired runtime change (\%)},legend style={at={(.5,1.03)},anchor=south,legend columns=2},ymajorgrids=true]
\addplot[ybar,bar width=\barwidth,fill=navy,draw=navy,error bars/.cd,y dir=both,y explicit] coordinates {(-0.17,-4.518702515943322) +- (0,0.013983612419588681) (0.83,-4.972759574002627) +- (0,0.01531263015113631) (1.83,-5.1080291189045575) +- (0,0.0006388687927925786)};\addlegendentry{E1 (3 pairs)}
\addplot[only marks,black,mark size=\pointsize,forget plot] coordinates {(-0.23,-4.530613890120525) (-0.17,-4.503305918400253) (-0.11000000000000001,-4.522187739309188) (0.77,-4.966771331058025) (0.83,-4.961345976988894) (0.8899999999999999,-4.990161413960962) (1.77,-5.107775364024343) (1.83,-5.107556113289411) (1.8900000000000001,-5.108755879399918)};
\addplot[ybar,bar width=\barwidth,fill=burnt,draw=burnt,error bars/.cd,y dir=both,y explicit] coordinates {(0.17,37.66690998422855) +- (0,8.54813252860203) (1.17,19.713397108644475) +- (0,4.078864128468423) (2.17,-6.42312058356128) +- (0,5.256522242165112)};\addlegendentry{E2 (5 pairs)}
\addplot[only marks,black,mark size=\pointsize,forget plot] coordinates {(0.11000000000000001,36.10902959910078) (0.14,44.07596239300247) (0.17,30.693839452395775) (0.2,48.675331167208185) (0.23,28.780387309435525) (1.1099999999999999,17.63669140825637) (1.14,18.886500387722457) (1.17,20.795064347728964) (1.2,15.205769972306939) (1.23,26.042959427207634) (2.11,-0.7610558601580898) (2.14,-6.752098565253035) (2.17,-2.690881749745551) (2.1999999999999997,-14.371718841826734) (2.23,-7.539847900822994)};
\addplot[black,no marks,forget plot] coordinates {(-.6,0) (2.6,0)};
\end{axis}\end{tikzpicture}
\par\raggedright\small
Bars: mean paired runtime change. Whiskers: SD. Dots: pairs. Negative values favor Overloaded.\par
\medskip
\begin{tabular}{llrr}
\toprule Group & Length & Mean change & ANOVA $p$\\\midrule
E1 & Short & $-4.52\%$ & $3.45\times10^{-6}$\\
E1 & Medium & $-4.97\%$ & $3.42\times10^{-6}$\\
E1 & Long & $-5.11\%$ & $7.66\times10^{-9}$\\
E2 & Short & $+37.67\%$ & 0.000475\\
E2 & Medium & $+19.71\%$ & 0.000371\\
E2 & Long & $-6.42\%$ & 0.056607\\\bottomrule
\end{tabular}\par
\medskip\normalsize
E1 is the M3 Max campaign ($n=3$, Arg1=360). E2 uses Jacob’s Intel/Ubuntu environment ($n=5$, Arg1=275/400/550). Workload parameters and repetition strategies differ, so these results cannot isolate an environment effect.
\smallsource{Updated v8 data, September 22. Welch between-group p: short .000383, medium .000173, long .6057.}
\end{block}
\vspace{.6cm}
\begin{block}{Sign consistency and mean differences are distinct}
All three E1 pairs improve at each length. In E2, all five pairs worsen at short/medium lengths and improve at long length.\par
\medskip
\small
\begin{tabular}{lrrr}
\toprule Group & $\chi^2$ & Approximate $p$ & Exact $p$\\\midrule
E1, each length & 3 & .0833 & .2500\\
E2, each length & 5 & .0253 & .0625\\\bottomrule
\end{tabular}\par
\medskip\normalsize
\textbf{No v8 directional test is significant at 0.05.} The exact test avoids misleading small-sample chi-square significance.
\end{block}
\vspace{.6cm}
\begin{block}{A resource improvement can accompany a slowdown}
In E2, minor faults decrease by approximately \textbf{77\%} at every size, while peak memory rises by \textbf{1.48--4.39\%}. Both effects are significant in the paired metric analyses.
\medskip
Short and medium runtime still increases. A lower fault count does not establish that allocation improves application runtime. All E2 major-fault counts are zero, so their percentage change and ANOVA are undefined.
\smallsource{Updated Metric ANOVA. Changes compare means. Resource p values do not establish causal contributions to runtime.}
\end{block}
\end{column}
\begin{column}{.30\paperwidth}
\begin{block}{CMA did not consistently worsen pthreads}
\centering\renewcommand{\plotheight}{17cm}\begin{tikzpicture}\begin{axis}[width=.94\linewidth,height=\plotheight,axis lines=left,tick label style={font=\footnotesize},label style={font=\small},legend style={font=\footnotesize,draw=none},ybar,xmin=-.6,xmax=3.6,ymin=0,ymax=460,xtick={0,1,2,3},xticklabels={WSL v4,WSL v6,Linux v4,Linux v6},ylabel={Mean runtime (s)},legend style={at={(.5,1.03)},anchor=south,legend columns=2},ymajorgrids=true]
\addplot[ybar,bar width=\barwidth,fill=navy,draw=navy,error bars/.cd,y dir=both,y explicit] coordinates {(0,395.795) +- (0,3.762) (1,292.769) +- (0,1.618) (2,348.573) +- (0,2.909) (3,224.896) +- (0,1.057)};\addlegendentry{Regular}
\addplot[ybar,bar width=\barwidth,fill=burnt,draw=burnt,error bars/.cd,y dir=both,y explicit] coordinates {(0,397.664) +- (0,6.023) (1,300.914) +- (0,.627) (2,246.981) +- (0,2.615) (3,226.776) +- (0,.843)};\addlegendentry{CMA v7}
\end{axis}\end{tikzpicture}
\par\raggedright\small
Same i7-11370H host, 16 GB RAM, eight threads, $n=800$. Means and SD from the report. Pthreads v4 uses queues/mutexes; v6 removes them. Both use CMA v7 or Regular new.\par
\medskip
\begin{tabular}{llrl}
\toprule Environment & Pthreads & CMA change & Tukey $p$\\\midrule
WSL & v4 & $+0.47\%$ & .8495\\
WSL & v6 & $+2.78\%$ & .0135\\
Linux & v4 & $-29.15\%$ & $<.0001$\\
Linux & v6 & $+0.84\%$ & .4965\\\bottomrule
\end{tabular}\par
\medskip\normalsize
Central locking, bin searches, and row placement can affect parallel execution. \textbf{Their individual costs were not isolated.} The report supports mixed outcomes, including a substantial Linux v4 improvement.
\smallsource{Report XII, Tables XII, XIV, XVI and XVIII. Pthreads versions and allocator versions identify different components.}
\end{block}
\vspace{.6cm}
\begin{block}{Breadth and limits of environmental variation}
Tests span \textbf{ARM64 and x86-64}, including M3/M5, Intel, and Ryzen systems. Environments include macOS, native Linux, and Windows WSL. Toolchains and OS versions differ.
\medskip
Coverage is not a balanced factorial experiment. Fixed-size repetition and problem-size growth answer different questions. Background activity, thermals, caches, and execution order can affect results.
\medskip
Three or five observations provide limited assumption checks. A large Shapiro--Wilk or Levene p value means no detected violation, not proof that the assumption holds.
\smallsource{Report V, VI-F and XIV; shared inventory and new v8 CSV. Some environment metadata remain missing.}
\end{block}
\vspace{.6cm}
\begin{block}{Conclusion}
\textbf{Allocator value depends on the request pattern and execution environment.} Representative workloads reveal tradeoffs. Effect sizes and inferential tests qualify claims. Code inspection suggests mechanisms; lock-wait and cache/TLB profiling would be needed to establish why a particular run wins or loses.
\end{block}
\vspace{.4cm}
{\footnotesize Sources: Group 2, \emph{Custom Memory Allocator Profiling}, supplied 31-page report; updated v8 analysis and September 22 CSV. Methods: NIST Sign Test; SciPy paired-test and Welch ANOVA documentation. The accompanying slides and source bundle provide detailed tables and source mappings.}
\end{column}
\end{columns}
\end{frame}
\end{document}
